\section{One actual finite polynomial family on the original real central domain}

Finite original homogeneous relations represent the ENTIRE original initial
ideal. The exact real central ideal equality allows their actual functions to
present every nearby WHOLE central stalk on one real open. Each coefficient is
an actual holomorphic function on its stated domain; no origin-germ evaluation
at arbitrary nearby points is introduced. All original inputs remain original
regular geometry and original coherence.

\begin{theorem}[Original input spreads the entire finite original relation model]
\label{hir:real-domain-original-entire-finite-central-relation-model}
\source{bierstone-milman08-functoriality}
\uses{hir:real-domain-original-whole-central-initial-ideal-equality,hir:relative-normal-deformation-finite-functional-model,hir:relative-normal-deformation-finite-central-spread,hir:real-domain-coherence-actual-finite-analytic-presentation,hir:real-domain-whole-origin-local-dictionary}
\lean{Hironaka.HironakaSmoothCentreRealDomainFinitePolynomialModel.exists_original_regular_real_central_finite_relation_model}
\leanok
Original regular geometry and coherence internally supply the genuine
flattening, selected index and real domains for the whole graph and central
family. On one smaller real central open containing zero and contained in
both the original real central domain and the genuine flattening target,
EVERY actual central membership function is exactly the finite analytic span
of the actual representative functions of the original finite homogeneous
relations. The same selected index and the whole all-point colon and
saturation-union dictionaries on the actual graph open are retained.
\end{theorem}
\begin{proof}
Use the original real whole central equality to supply the stated actual
flattening, selected index, graph open and central domain. The finite original
homogeneous generators give their actual holomorphic representative functions
on one open containing zero; their finite function span has genuine local
finite generation at every point of that open. At zero its full representative
dictionary is precisely the whole analytic initial-image ideal. The actual
central domain coherence supplies genuine finite analytic generation at zero,
so the local whole origin dictionary and the supplied whole central equality
identify EVERY central representative with that same finite relation span.
Two genuinely finitely presented analytic ideals equal at zero agree on a
real open: express each finite generator of each presentation in the other,
and intersect the finitely many neighborhoods where these actual equalities
and coefficient holomorphicity hold. Their finite spans then agree at every
point of that intersection. Intersect this equality open with the actual
finite-function open, real central domain and genuine flattening target.
The resulting open contains zero and retains the full finite span dictionary
at every point. Restrict central coherence there and retain the actual graph
open and all its exact colon and union identities.
\end{proof}

\begin{theorem}[The same real central family has actual homogeneous polynomial functions]
\label{hir:real-domain-original-whole-central-polynomial-functions}
\source{bierstone-milman08-functoriality}
\uses{hir:real-domain-original-entire-finite-central-relation-model,hir:relative-normal-central-literal-polynomial-functions,hir:relative-normal-central-finite-polynomial-representative-comparison,hir:relative-normal-analytic-coefficients-whole-polynomial-dictionary,hir:relative-normal-graded-whole-initial-ideal}
\lean{Hironaka.HironakaSmoothCentreRealDomainFinitePolynomialModel.exists_original_regular_real_central_polynomial_model}
\leanok
At the SAME selected index the existing central ideal has a complete actual
finite polynomial-function presentation on one real open containing zero
inside its true central domain and the genuine flattening target. Its finite
family is the original homogeneous relation family. Every actual coefficient
representative is holomorphic on its real base domain; all polynomial
functions are holomorphic there, and every original relation retains its
actual homogeneous degree under all specializations on that open. The actual
finite polynomial-function span presents every WHOLE central stalk.
\end{theorem}
\begin{proof}
Take the actual finite central relation model from the preceding theorem.
For each of its finitely many original polynomials, the genuine whole germ of
its literal polynomial function is its entire analytic polynomial evaluation,
which is also the whole germ of its chosen actual relation representative.
Their actual functions agree near zero. Intersect the finitely many equality
opens, so their complete finite spans agree at every point on that open.
The finitely many coefficient germs have actual holomorphic representatives
on a common real base open containing zero; zero coefficients use the actual
zero representative, so the common holomorphicity statement covers all
exponents. Pull that open back by the true tangent projection and intersect
with the preceding real central open and finite-function equality open.
Finite products and sums of actual holomorphic coefficient pullbacks and
normal coordinate monomials make each literal polynomial function holomorphic
at every point of this intersection. Their finite span therefore has local
finite analytic generation there. The original finite family consists of
actual homogeneous original relations, and evaluation of its actual
coefficient representatives preserves their homogeneous degree. Combining
the two finite-span dictionaries gives the entire central membership
presentation at every point. All retained actual graph saturation identities
still use the same original index and graph open.
\end{proof}

\begin{theorem}[One real base open gives all entire translated central polynomial kernels]
\label{hir:real-domain-original-all-translated-central-polynomial-kernels}
\source{bierstone-milman08-functoriality}
\uses{hir:real-domain-original-whole-central-polynomial-functions,hir:relative-normal-nearby-whole-polynomial-evaluation-dictionary,hir:relative-normal-nearby-entire-polynomial-coefficient-germs,hir:relative-normal-homogeneous-whole-analytic-contraction,hir:real-centre-domain-whole-nearby-translated-ideals}
\lean{Hironaka.HironakaSmoothCentreRealDomainFinitePolynomialModel.exists_original_regular_real_translated_central_polynomial_kernels}
\leanok
Original inputs supply one real base open containing zero such that every
actual base point maps into the genuine flattening target and coordinate
centre plane. At every point of this SAME open, the WHOLE origin ideal of the
actual translated central family is the analytic evaluation image of the
homogeneous ideal generated by the same finite original polynomials with
actual translated coefficient germs. Its ENTIRE polynomial preimage is that
same homogeneous ideal. Every translation uses actual functions holomorphic
on its stated real domain; no global model coherence is an input or conclusion.
\end{theorem}
\begin{proof}
Take the real central polynomial open $V$ supplied above and let the base open
be the true tangent-inclusion preimage of $V$. It is open and contains zero;
its actual inclusion lies in the genuine flattening target because
$V$ does, and it lies in the coordinate centre plane by the inclusion's
literal coordinate identities. Each actual coefficient representative is
holomorphic near every base point by the common real coefficient condition.
Translate these actual functions to zero to obtain genuine nearby coefficient
germs. The true polynomial translation formula identifies every translated
literal polynomial-function germ with the analytic evaluation of the same
original polynomial with these actual nearby coefficients. The full local
finite-span pullback dictionary presents the translated central family by
all those actual translated functions. Passing its complete finite equality
to germs in both directions identifies its ENTIRE origin ideal with their
whole ideal span. The genuine translation polynomial formulas identify that
span with the analytic image of the displayed finite polynomial ideal.
Every displayed polynomial retains its original homogeneous degree, so their
ideal is homogeneous. Genuine whole homogeneous analytic contraction then
makes the entire polynomial preimage of its analytic ideal image exactly the
same homogeneous ideal. All points use one actual base open and one finite
original relation family, including empty families and zero normal dimension.
This is the exact kernel for the translated ACTUAL central ideal; the
identification with each original nearby initial ideal and subsequent common
free open remain separate obligations.
\end{proof}
